\section{Dyadic critical-source localization at the original roof}
\label{sec:v82-critical-dyadic}
This section isolates a quantitative endpoint criterion without changing the
physical source, the exact displacement and occupation coefficient, or the
terminal membership test. All statements below are at fixed collision count
until the order of limits is specified. The purpose of the decomposition is
to avoid exchanging fixed-width concentration with an unbounded Abel kernel.

Let $u=u_{n,k,m,R}$ be the nonnegative complete inverse-incidence density
from \eqref{eq:v80-depth-insertion}, and let
$v^\eta\,dt=L_\#(I_m|g|^{-\eta}\one_\Pi\nu_R^*)$ for $0<\eta<1/2$.
The inherited common coarea kernels give measurable nonnegative versions.
Put $I_j(t)=(t+2^{-j-1}r,t+2^{-j}r]$ for $j\ge0$.

\begin{lemma}[Exact dyadic majorant for the critical Abel source]
\label{lem:v82-dyadic}
For any nonnegative locally integrable $v$, any $0<r<\infty$,
and almost every $t$,
\begin{equation}\label{eq:v82-dyadic}
 \int_0^r s^{-1/2}v(t+s)\,ds
 \le \sqrt2\sum_{j=0}^\infty
 (2^{-j-1}r)^{-1/2}
 \int_{I_j(t)}v(t+s)\,ds .
\end{equation}
The inequality holds in $[0,\infty]$, with no regularity assumption on
$v$. In particular, if for each fixed outer parameter $\varepsilon$
one proves a collision-uniform *local* bound
\begin{equation}\label{eq:v82-hypothesis}
 \limsup_m\sup_{R,n,k}
 m^2\operatorname*{ess\,sup}_{t\in\mathcal C(M)}
 \int_0^\ell v(t+s)\,ds\le H_{\varepsilon,M}(\ell)
 \quad (0<\ell<r_0),
\end{equation}
and the limsup is uniform over the finite partial sums in
\eqref{eq:v82-dyadic} including their tails, then the ordered Abel
height is bounded by
\begin{equation}\label{eq:v82-criterion}
 C\sum_{j\ge0}(2^{-j}r)^{-1/2}
 H_{\varepsilon,M}(2^{-j}r).
\end{equation}
For $H_{\varepsilon,M}(\ell)\le
A_M\varepsilon^{-q}\ell^{1/2+\delta}$, $\delta>0$, this is
$O_M(\varepsilon^{-q}r^\delta)$. The tail-uniformity clause cannot
be omitted: pointwise limsups of every fixed scale do not control
the limsup of an infinite sum of nonnegative terms.
\end{lemma}
\begin{proof}
Partition $(0,r]$ into the half-open intervals
$(2^{-j-1}r,2^{-j}r]$. On each interval
$s^{-1/2}\le(2^{-j-1}r)^{-1/2}$; Tonelli proves
\eqref{eq:v82-dyadic}. Each interval integral is bounded by the
cumulative positive integral up to $2^{-j}r$.
For any finite partial sum the limsup is bounded by the sum of
the finitely many corresponding limsups. A uniform bound on the
remaining tails permits passage from these partial sums to the
infinite sum, with no exchange of collision limit and an unbounded
kernel. The geometric series with ratio $2^{-\delta}$ proves the
last estimate.
\end{proof}

\begin{proposition}[The exact-label surface obstruction remains independent]
\label{prop:v82-surface-independent}
The volume criterion \eqref{eq:v82-hypothesis}, even if established
with its tail uniformity, does not by itself bound the nongrazing
first-exit surface height $\mathfrak F_M(a,r)$ in
\eqref{eq:v80-endpoint-budget}. Its first-exit selection,
one-sided physical components, backward lifetime and exact section
decisions remain part of a separate positive measure. Any
surface local theorem used in the endpoint budget must be
established for this actual measure, not inferred from bounded
volume multipliers or from cancellation of internal section currents.
\end{proposition}
\begin{proof}
The coarea source entering \eqref{eq:v82-hypothesis} is a
two-dimensional volume pushforward. The outgoing measure in
\eqref{eq:v80-outgoing-assignment} is a one-dimensional transverse
flux, supported on true boundary faces, with first-exit selection.
An estimate for the former supplies no absolute-continuity or
uniform trace estimate for the latter. The assertion is a precise
separation of hypotheses, not a claim of new decay for the surface
measure.
\end{proof}
